\chapter{\MakeUppercase{Project Implementation}}
\section{Implementation Overview}
The project is implemented in Python. The file \verb|environment.py| defines the disaster \ac{dtn} testbed, simulation constants, node creation, message generation, mobility update, contact detection, delivery checking, expiration, and metric computation. The file \verb|message.py| defines message identity, source, destination, creation time, critical flag, hop count, and copy count. The file \verb|node.py| defines node identity, role, mobility properties, location, and buffer.

The routing implementations are placed inside the \verb|routing| package. The baselines are implemented in \verb|routing/epidemic.py| and \verb|routing/spray.py|. The proposed protocol is implemented in \verb|routing/spatio_semantic.py|. The experiment runner in \verb|main.py| executes every protocol across all traffic levels and writes aggregate outputs to the \verb|plots| directory.

\section{Environment Implementation}
The environment creates four node classes:

\begin{itemize}
	\item 35 civilians with slow movement.
	\item 8 responders with medium movement.
	\item 3 static shelters clustered in one region.
	\item 5 drones with high movement speed.
\end{itemize}

At every time step, the environment probabilistically generates messages, moves nodes, detects contacts, invokes router exchanges, checks deliveries, and removes expired messages. This implementation gives all protocols the same external conditions and isolates differences to their routing decisions.

\section{Routing Baseline Implementation}
The Epidemic router clones messages to any encountered node that lacks the message. It increments the transmission counter for every successful transfer and increments the drop counter when the receiver buffer is full.

The Spray and Wait router uses message copy counts. When a sender has more than one copy, it gives one copy to a receiver that does not already hold the message. When a message has only one copy, it is forwarded only to the destination. This design reduces unnecessary replication but can slow delivery.

\section{Spatio-Semantic Router Implementation}
The proposed router maintains two tracking tables on nodes:

\begin{itemize}
	\item Encounter table: records decayed contact frequency with other nodes.
	\item Zone table: records decayed visit frequency for spatial zones.
\end{itemize}

The router updates these tables during exchanges. It then computes the utility of the sender and receiver relative to the message destination. If the receiver is the destination, delivery is always attempted. Otherwise, the protocol uses message priority and copy count to decide whether to spray, utility-forward, or wait.

\section{Critical Buffer Reserve}
The implementation reserves 35 percent of the buffer for critical messages. Non-critical messages can only occupy the non-reserved portion. Critical messages can use any available slot and can evict a non-critical message if necessary. This prevents non-critical traffic from consuming all storage during high-load conditions.

\section{Result and Visualization Implementation}
The file \verb|dashboard.py| loads the three result files:

\begin{itemize}
	\item \verb|plots/epidemic_results.csv|
	\item \verb|plots/spray_results.csv|
	\item \verb|plots/spatio_semantic_results.csv|
\end{itemize}

The dashboard provides trend graphs, radar charts, raw data tables, and a launch button for the visual simulation. The file \verb|modern_multi_demo.py| displays a full-screen three-panel comparison of Epidemic, Spray and Wait, and Spatio-Semantic routing.

\begin{figure}[h!]
    \centering
    \fbox{\parbox{0.88\textwidth}{\centering
    Image placeholder: add a screenshot of the Raw Data tab from the dashboard after all experiments have been generated.}}
    \caption{Raw aggregated simulation results in the dashboard}
    \label{fig:raw_data_placeholder}
\end{figure}
